\chapter{Introducción}\label{chap:introduccion}

\section{Motivación}
ROS2 (\textit{Robot Operating System 2}) es un conjunto de herramientas de \textit{software} 
que proporcionan una plataforma estandarizada para desarrollar aplicaciones 
de robótica \parencite{rosROSROS}.
Esta tecnología se caracteriza por ser una solución multidominio (permite
desarrollar aplicaciones para robots externos e internos, industriales 
o domésticos, terrestres, aereos o marítimos, etc.) y 
multiplataforma (es soportado por diferentes sistemas operativos 
Linux, Windows y MacOS además de otras plataformas embebidas). 
Por otro lado, se trata de un conjunto de herramientas de código abierto
(\textit{open source}) distribuida a través de licencias permisivas 
(como Apache 2) que facilitan la distribución de soluciones 
personalizadas propietarias \parencite{rosROSROS}.
La adopción de las soluciones proporcionadas por ROS2 en el año 
2023 experimentó un incremento del 22,9\% de diferentes compañías \parencite{ros2Metrics}, lo que 
hace de esta tecnología un objetivo claro de investigación.

La interfaz proporcionada por ROS2 permite desarrollar nuevos métodos de 
interacción humano-robot de forma que sea posible visualizar distintos 
parámetros de la situación actual del sistema o teleoperar los diferentes 
dispositivos desplegados. Una de las maneras más habituales de realizar 
las operaciones aneteriormente mencionadas es empleando un computador 
personal (PC) ya sea o no portátil, lo que limita la movilidad y maniobrabilidad 
del operador del robot.

Existen tendencias más actuales que emplean dispositivos montados en la 
cabeza (HMD; \textit{Head Mounted Displays}) como \textit{Microsoft HoloLens},
\textit{MetaQuest} o \textit{HTC Vive} que permiten mezclar el mundo 
real y el virtual mediante la realidad aumentada empleando interfaces más 
intuitivas, directas y que permiten la movilidad total o parcial del operador.
Esto facilita el aprendizaje, la operativa y la solución de errores a la vez que 
reduce la carga cognitiva necesaria con respecto a la utilización de interfaces 
gráficas bidimensionales no contextualizadas \parencite{richardsetal, robotics14080106}.
Sin embargo, estos dispositivos continuan teniendo un coste demasiado elevado
para muchos usuarios, lo que 
limita la capacidad de despliegue de las soluciones 
desarrolladas a las organizaciones más pudientes. 
Además, requieren distribuir el centro operativo 
en diferentes instrumentos en lugar de centralizar las 
acciones y datos en un único dispositivo como 
los \textit{smartphones} o \textit{tablets} que cumplen múltiples 
propósitos.

A pesar de las aparentes ventajas económicas y operacionales de 
emplear los dispositivos móviles como interfaz en la 
visualización de datos o tele-operación de robots, no existen 
demasiadas soluciones integradas que permitan realizan dichas 
acciones empleando realidad aumentada.


\section{Propuesta y objetivos}
El proyecto propone desarrollar una interfaz multiplataforma (Android/iOS) 
que permita la interacción con robots (visualización de datos y 
tele-operación) basados en ROS2 
mediante realidad aumentada eliminando las barreras materiales 
y económicas existentes en
la teleoperación y visualización de datos.

Para ello, se deben evaluar las soluciones existentes y su compatibilidad
con ROS2 y la realidad aumentada móvil para adoptarlas y adaptarlas en caso de 
ser necesario. 
También es necesario crear un sistema de comunicaciones ligero
y multiplataforma coordinado preferiblemente desde el dispositivo móvil eliminando
la necesidad de agentes coordinadores externos.

Para proporcionar la visualización de datos y la teleoperación es preciso 
incorporar interfaces en realidad aumentada intuitivas de manera que los 
operadores puedan contextualizar en todo momento la situación del sistema 
y tomar las medidas necesarias en el desarrollo de su trabajo (i.e. cambiar
parámetros de los agentes).

Es necesario, además, realizar una evaluación de la usabilidad y la latencia 
del sistema de manera que se cuente con una fiabilidad suficiente como 
para poder desplegar la solución en entornos de operación reales.

\section{Estructura del documento}
El documento se estructura en diferentes secciones en las que se realizará un recorrido de la situación actual de la oferta 
de aplicaciones de visualización de parámetros y variables de robots en diferentes plataformas haciendo incapié en aquellas
que proporcionan movilidad y ubicuidad a los usuarios como son los dispositivos móviles. Esta investigación se englobará en la
sección de estado del arte (Capítulo \ref{chapter:estado_del_arte}).

Tras vislumbrar la posible necesidad de nuevas propuestas en este ámbito, se presentarán los materiales y métodos (Capítulo \ref{chapter:materiales_metodos})
necesarios para llevar a cabo un desarrollo exitoso de una plataforma de realidad virtual en la que se puedan realizar 
las funciones básicas presentes en aplicaciones de esta índole (como la publicación de mensajes, la visualización de mensajes, 
el estado de las partes del robot, sus velocidades, etc.).

Los resultados (Capítulo \ref{chapter:resultados}) obtenidos serán presentados mostrando visualmente el \textit{software} obtenido tras el proceso de desarrollo, 
repasando los objetivos y su estado de cumplimiento y las pruebas de integración realizadas para validar el uso básico de la aplicación.
En este apartado se presentarán también ciertas limitaciones encontradas durante el desarrollo que llevaron a rediseños de la planificación 
original.

Finalmente, se presentarán las conclusiones (Capítulo \ref{chapter:conclusiones}) de la realización del proyecto y el aprendizaje asociado a las diferentes tecnologías mediante las 
líneas de trabajo futuras.

Como elementos adicionales, se incluyen sendos manuales de instalación (Apéndice \ref{chapter:manual_instalacion}) y de uso (Apéndice \ref{chapter:manual_uso}), 
de forma que se guíe al usuario desde la intencionalidad 
de emplear la aplicación hasta su utilización real.

